% ---------------------------------------------------------------------
%  Mode d'emploi de l'ouvrage
% ---------------------------------------------------------------------
\chapter*{Comment utiliser cet ouvrage}
\addcontentsline{toc}{chapter}{Comment utiliser cet ouvrage}
\markboth{Comment utiliser cet ouvrage}{Comment utiliser cet ouvrage}

\section*{Organisation}

L'ouvrage comprend deux parties, suivies d'annexes.

La \textbf{partie I} pose les fondements. Le chapitre~\ref{ch:fondements} présente l'architecture
Lakehouse, ses origines et ses principes : formats de fichiers et de tables ouverts, séparation du
stockage et du calcul, architecture en couches Bronze, Silver et Gold, modélisation dimensionnelle,
traitement par lots et en continu. Le chapitre~\ref{ch:environnement} installe l'environnement de
travail utilisé dans toute la suite.

La \textbf{partie II} est formée de quatre travaux pratiques (TP), qui construisent progressivement
une plateforme complète sur les données d'un site de commerce en ligne.

\begin{center}
\small
\renewcommand{\arraystretch}{1.3}
\begin{tabularx}{\textwidth}{@{}l X l@{}}
\toprule
\textbf{Chapitre} & \textbf{Thème} & \textbf{Durée indicative} \\
\midrule
\ref{ch:tp1}, TP 1 & Données structurées : de la base de vente au tableau de bord, Iceberg et Delta Lake & 6 h + 2 h 30 \\
\ref{ch:tp2}, TP 2 & Texte et JSON : avis clients, suivi des colis, modélisation avec dbt & 6 h + 2 h \\
\ref{ch:tp3}, TP 3 & Images, orchestration et BI : Airflow, Trino, Superset, supervision & 9 h + 3 h \\
\ref{ch:tp4}, TP 4 & Batch ou streaming : Kafka et Spark Structured Streaming & 6 h + 2 h \\
\bottomrule
\end{tabularx}
\end{center}

La durée indique le temps en salle, puis le travail personnel conseillé. Les TP se suivent : chacun
utilise les tables construites par les précédents. Les annexes regroupent les aide-mémoires de code,
un glossaire français-anglais et la description des jeux de données avec leurs licences.

\section*{Le matériel d'accompagnement}

Les énoncés des TP prennent la forme de \emph{notebooks} Jupyter, accompagnés de scripts et d'un
environnement Docker complet, disponibles à l'adresse \depotcode. Le chapitre~\ref{ch:environnement}
explique comment les installer. Le lecteur copie les énoncés dans son dossier de travail et complète
les cellules marquées \cmd{À COMPLÉTER}. Chaque chapitre de la partie II indique le notebook à
utiliser pour chaque partie, et les numéros de cellule à exécuter avant chaque capture d'écran.

Les résultats chiffrés cités dans les chapitres (nombre de lignes, de doublons, de fichiers\dots)
ont été obtenus avec les versions des données et des logiciels fournies. Le lecteur doit retrouver
exactement les mêmes nombres. Seules les durées varient d'une machine à l'autre.

\section*{Conventions}

\begin{center}
\small
\renewcommand{\arraystretch}{1.3}
\begin{tabularx}{\textwidth}{@{}l X@{}}
\toprule
\textbf{Élément} & \textbf{Signification} \\
\midrule
\textbf{\textcolor{lhblue}{Notion clé :}} & une notion à retenir, définie au moment où elle sert \\
\textbf{\textcolor{lhred}{Attention :}} & un piège fréquent, ou une précaution à prendre \\
\textbf{Résultat attendu :} & ce que le lecteur doit obtenir à la fin d'une partie \\
\etape{1}\ \etape{2}\dots & les étapes numérotées des figures d'architecture \\
\nb{01\_bronze} & le nom d'un notebook \\
\cmd{silver.orders} & un nom de table, de fichier, de colonne ou une commande \\
\bottomrule
\end{tabularx}
\end{center}

Les commandes à saisir dans un terminal sont présentées sur fond sombre, le code Python ou SQL sur
fond clair :
\begin{lstlisting}[style=shell]
docker compose up -d
\end{lstlisting}
\begin{lstlisting}[language=Python]
spark.table("lakehouse.silver.orders").count()
\end{lstlisting}

Les termes techniques sont donnés en français, avec leur équivalent anglais à la première
occurrence lorsque celui-ci est d'usage courant dans la profession. Le glossaire de
l'annexe~\ref{ann:glossaire} les rassemble.
